% ComplaintLens: technical report (Deliverable 4). Every number cites the results file in
% reports/results/ that produced it. Items marked \todo{} are pending experiments: they must be
% filled from real runs, never estimated.
\documentclass[11pt,a4paper]{article}
\usepackage[margin=2.3cm]{geometry}
\usepackage[T1]{fontenc}
\usepackage[utf8]{inputenc}
\usepackage{lmodern,microtype,booktabs,tabularx,xcolor,hyperref,graphicx}
\usepackage[numbers,sort]{natbib}
\hypersetup{colorlinks=true,linkcolor=black,citecolor=blue!50!black,urlcolor=blue!50!black}
\newcommand{\todo}[1]{\textcolor{red}{[TODO: #1]}}
\newcommand{\res}[1]{\texttt{\small reports/results/#1}}

\title{ComplaintLens: Bilingual Complaint Understanding, Routing and Verified Response\\
with Guaranteed-Risk Automation}
\author{Mohamed Aziz Mansour \and Rimel Zouari\\[2pt]
\small NLP Semester Project 2026--2027 --- Dr.\ Nizar Omheni}
\date{\todo{submission date}}

\begin{document}
\maketitle

\begin{abstract}
Customer-service teams receive complaints through many channels and languages and must decide, for
each message, whether it needs an answer, what it is about, whether it is dangerous, and what to
reply. We present ComplaintLens, an English/French system that turns one incoming message into a
structured, traced decision: message type, category, severity and risk flags, the linked customer
record, a reply grounded in company policy and in facts computed from the record, a deterministic
verification of that reply, and an automation decision with a statistical error bound. We train and
evaluate on real data (expert-labelled complaint tweets, 41{,}855 real Uber support conversations)
and deploy the system as the live support desk of a fictional ride and food-delivery company.
A TF-IDF baseline already exceeds every published non-neural result on complaint identification
(macro-F1 81.5 vs.\ 79.0); dense multilingual retrieval beats BM25 on customer-worded queries
(MRR 0.61 vs.\ 0.49, $p=0.022$) while equal-weight rank fusion does not; a rate-based burst detector
recovers 94\% of injected 20-message bursts where a count threshold recovers none, and flags a real
app outage one hour after onset; and a Learn-then-Test gate automates 37\% of messages at a
guaranteed 10\% error bound, refusing to automate when the evidence cannot support a stricter bound.
\todo{headline numbers for the fine-tuned encoder (RQ1) and grounded generation (RQ3)}
\end{abstract}

\section{Introduction and problem statement}
\label{sec:intro}
\paragraph{Task.} Input: one customer message (text, channel, sender, thread). Output: language
(EN/FR/other); message type (complaint, request, feedback, praise, spam) and whether a response is
required; category (11 classes); expressed severity (4 levels, following \citet{trosborg1995pragmatics}
as operationalised by \citet{jin2021severity}); risk flags (safety, injury, harassment, legal, fraud,
media, vulnerable person); entities linked to the customer's trip or order; an explainable priority
and team; a reply grounded in cited policy and record facts; and an automation decision
(send / await approval / escalate / no reply).

\paragraph{Why it matters.} Answering late or wrongly is costly; answering dangerous messages
automatically is unacceptable. Our design therefore separates \emph{how strongly} a complaint is
expressed (severity) from \emph{how risky} it is (risk flags), and never automates safety, injury,
harassment or legal cases.

\paragraph{Research questions.}
\textbf{RQ1} Does a small fine-tuned multilingual encoder, distilled from LLM silver labels and human
labels, match its LLM teacher and beat classical baselines, and how well does it transfer to French?
\textbf{RQ2} Can the reply decision be automated with a guaranteed error bound, and at what
automation rate?
\textbf{RQ3} Do retrieval, record facts and a claim verifier make replies more faithful, and how do
they compare with real agent replies?
\textbf{RQ4} Does burst detection find emerging problems earlier than count thresholds?

\section{Related work}
Complaint identification on social media was introduced by \citet{preotiuc2019complaints}, who
released 3{,}449 expert-labelled tweets; \citet{jin2021severity} added four severity levels and
reported transformer results (macro-F1 up to 88.2 for complaints and 55.7 for severity).
\citet{hardalov2018support} studied response generation on the Customer Support on Twitter corpus.
Retrieval combines BM25 \citep{robertson2009bm25} with multilingual sentence embeddings
\citep{reimers2020multilingual}, fused by Reciprocal Rank Fusion \citep{cormack2009rrf}; generation
grounded on retrieved text follows \citet{lewis2020rag}. Burst detection in streams goes back to
\citet{kleinberg2003bursty}. Our automation gate follows Learn-then-Test risk control
\citep{angelopoulos2021ltt}. Significance testing follows \citet{dror2018hitchhiker} and
\citet{bergkirkpatrick2012significance}.

\section{Data}
\label{sec:data}
\begin{table}[h]
\centering\small
\begin{tabularx}{\linewidth}{lXll}
\toprule
Dataset & Content & Size & Licence \\
\midrule
Customer Support on Twitter & Real customer/company tweets, 108 brands, with timestamps and agent replies & 2.81M tweets & CC BY-NC-SA 4.0 \\
\quad Uber subset & First contacts to @Uber\_Support, Oct--Dec 2017 & 41{,}855 conversations & (same) \\
Complaints on Twitter & Complaint vs.\ not, 9 domains \citep{preotiuc2019complaints} & 3{,}449 tweets & research use \\
\quad severity & 4 levels on the 1{,}232 complaints \citep{jin2021severity} & 1{,}232 & research use \\
LunaGo Help Center & 9 bilingual policy articles (ours, fictional) & 35 sections $\times$ 2 & ours \\
LunaGo records & Synthetic riders, trips, orders, charges (demo only) & generated & ours \\
Gold sets & Human annotation, EN + FR \todo{sizes, $\kappa$ per field} & \todo{} & ours \\
\bottomrule
\end{tabularx}
\caption{Datasets. Exploratory statistics: \res{eda\_uber\_first\_contacts.json}.}
\end{table}

\paragraph{What exploration changed in the design.}
(i) 95\% of Uber first contacts are one tweet (median 22 words): short-text methods matter.
(ii) About 49\% mention rides and 16\% food delivery, so the fictional company offers both.
(iii) The corpus was collected around company replies: only 19 of 19{,}762 customer-started threads
tagging @Uber\_Support went unanswered, so ``did the company reply?'' cannot label ``response
required''; that label comes only from human annotation.
(iv) The busiest day, 14 Nov 2017 (1{,}333 conversations), is a real incident: a free-doughnut
promotion that crashed the food-delivery app, our real-event test for RQ4.
(v) The severity release cannot be joined by ID (IDs corrupted to floats such as \texttt{2.26E+16})
nor by exact text (lower-cased, re-tokenised, \texttt{<user>}/\texttt{<url>} placeholders, one row
replaced by the spreadsheet error \texttt{\#NAME?}). Both releases share row order and complaint
labels on all 3{,}449 rows; we align by position and verify each pair by text similarity
($\geq 0.9$ except the flagged row), keeping the original-cased text.

\paragraph{Preprocessing, each step justified.} No lower-casing, stop-word removal or accent
stripping for neural models (negation and casing carry meaning; pretrained tokenisers are used as
is). Brand names and handles are masked (\texttt{[COMPANY]}, \texttt{[USER]}) so that models learn
the complaint, not the brand. Consecutive tweets of one customer are merged. Email quotes and
signatures are removed. Personal data is masked before storage and before any external API call.
Duplicates are grouped so that they never cross train/test boundaries.

\section{System and methods}
\label{sec:method}
Every message runs through six traced stages; each stage records its output, confidence and time,
which the console displays live (the ``X-ray'').
\begin{enumerate}
\item \textbf{Normalise}: email quote and signature removal, Unicode normalisation.
\item \textbf{Mask personal data}: emails, phone numbers (North American and French formats) and
card numbers validated with the Luhn checksum, so amounts and order numbers are not masked.
\item \textbf{Language}: py3langid \citep{lui2012langid} restricted to EN/FR with a short-text rule;
chosen over a 170\,MB alternative for footprint (\todo{accuracy by text length vs.\ alternatives}).
\item \textbf{Understand}: three signals combined transparently: classical models (TF-IDF word
1--2 + character 2--5 grams, logistic regression), an LLM prompted with the \emph{same annotation
guidelines as the human annotators}, and a bilingual recall-first risk lexicon that can add, never
remove, safety flags. Disagreement between LLM and classical model is an uncertainty signal.
\item \textbf{Link}: the customer by sender identity; the trip by a transparent evidence score
(trip number +10, matching amount +4, date +3, problem visible in the record +3). The trip kind is
only a tie-breaker; without specific evidence the system asks instead of guessing. Facts about the
record are computed by code, never by the LLM.
\item \textbf{Respond}: retrieval over the Help Center; an LLM draft restricted to retrieved
sections and record facts, with mandatory citations; a deterministic verifier (every citation must
exist and every number in the reply must appear in a cited source); and the gate (Section~\ref{sec:gate}).
Safety, injury, harassment and legal cases are escalated with a fixed safe acknowledgement; money is
never granted automatically.
\end{enumerate}

\section{Experiments and results}
All numbers come from scripts in \texttt{ml/scripts/}; seeds and configurations are stored in the
results files.

\subsection{Complaint identification and severity (baselines; RQ1 pending)}
Protocol: stratified 10-fold cross-validation with duplicates grouped, repeated over 5 seeds; default
hyperparameters; macro-F1 (classes are imbalanced).

\begin{table}[h]
\centering\small
\begin{tabular}{lcc}
\toprule
Model & Complaint vs.\ not & Severity (1--4) \\
\midrule
Majority (ours / published) & 39.1 / 39.1 & 13.1 / 13.0 \\
Best published non-neural & 79.0 & 43.5 \\
TF-IDF word + LR (ours) & 78.3 $\pm$ 0.2 & 44.6 $\pm$ 0.5 \\
TF-IDF char + LR (ours) & 79.7 $\pm$ 0.1 & 44.7 $\pm$ 0.6 \\
\textbf{TF-IDF word+char + LR (ours)} & \textbf{81.5 $\pm$ 0.2} & \textbf{46.2 $\pm$ 1.1} \\
Best published transformer & 88.2 & 55.7 \\
Fine-tuned multilingual encoder (ours) & \todo{} & \todo{} \\
\bottomrule
\end{tabular}
\caption{Macro-F1 (\%). Sources: \res{baseline\_complaints2019.json},
\res{published\_reference\_complaints2019.json}. Word+char beats char-only by +1.6 ($p=0.0005$) and
+2.6 ($p=0.013$), paired bootstrap.}
\end{table}

Our majority baselines reproduce the published values, an independent check of the data alignment.
Caveats: same protocol as the papers but different fold assignment, and no inner-loop tuning (which
favours the papers). Grouping duplicates changed scores by less than the seed variance
(\res{baseline\_complaints2019\_ungrouped\_ablation.json}).

\subsection{Retrieval (RQ3, retrieval part)}
\begin{table}[h]
\centering\small
\begin{tabular}{lcccc}
\toprule
Retriever & R@1 & R@3 & MRR@10 [95\% CI] & Latency \\
\midrule
BM25 (implemented from scratch) & 0.40 & 0.53 & 0.49 [0.39, 0.59] & 0.04 ms \\
\textbf{Dense, multilingual MiniLM} & \textbf{0.49} & \textbf{0.70} & \textbf{0.61 [0.52, 0.71]} & 4 ms \\
Hybrid BM25 + dense (RRF) & 0.40 & 0.63 & 0.54 [0.44, 0.63] & 4 ms \\
\bottomrule
\end{tabular}
\caption{70 customer-worded questions (35 EN, 35 FR). Dense vs.\ BM25: $+0.121$ MRR, $p=0.022$;
hybrid vs.\ BM25: $+0.044$, $p=0.06$. Source: \res{retrieval\_helpcenter.json}.}
\end{table}
Equal-weight fusion \emph{hurt} because BM25 is much weaker on customer wording. Searching with the
LLM's one-line English summary as well (sections share IDs across languages) raised MRR to 0.66
(\res{retrieval\_summary.json}). To fit the free host's 1\,GB container we deploy the model's
official int8 file: MRR 0.59 vs.\ 0.61 on the text alone (paired bootstrap $p=0.13$), same top
section for 89\% of questions, 0.64 with the summary (\res{retrieval\_int8.json}). This set was
used to choose the retriever, so it is a development set; \todo{final numbers on the held-out
question set written by the team; reranker; larger embedding models}.

\subsection{Emerging-issue detection (RQ4)}
\begin{table}[h]
\centering\small
\begin{tabular}{lcccc}
\toprule
Method & +10 msgs & +20 & +40 & Incidents/day (Oct / Nov) \\
\midrule
\textbf{Rate-based bursting terms (Poisson $z$)} & \textbf{0.73} & \textbf{0.94} & \textbf{1.00} & 2.08 / 4.94 \\
Fixed count threshold & 0.00 & 0.00 & 0.00 & 1.12 / 2.19 \\
\bottomrule
\end{tabular}
\caption{Recall on injected bursts of real messages (3 seeds $\times$ 60 trials); thresholds chosen on
October under an equal alarm budget ($\leq 3$ incidents/day) and applied unchanged to November.
Source: \res{trends\_uber.json}.}
\end{table}
On 14 Nov 2017 the top incident is ``krispy, kreme'' (102 messages, $z=101$); hour by hour, the
problem starts at 13:00 UTC and the first alert fires at 14:00, one hour after onset. November
exceeds the alarm budget partly because of genuine incidents that cannot be told apart from false
alarms without labels.

\subsection{Problem discovery}
\begin{table}[h]
\centering\small
\begin{tabular}{lccc}
\toprule
Method ($k=30$, 3 seeds) & NPMI coherence & Diversity & Largest cluster \\
\midrule
TF-IDF + $k$-means & 0.097 $\pm$ 0.011 & 0.78 & 18.6\% \\
LDA & 0.000 $\pm$ 0.010 & 0.95 & 8.7\% \\
Multilingual embeddings + $k$-means & 0.088 $\pm$ 0.009 & 0.78 & 5.4\% \\
\bottomrule
\end{tabular}
\caption{41{,}690 Uber first contacts; every method's clusters are described with the same c-TF-IDF
function. Source: \res{topics\_uber.json}.}
\end{table}
Coherence does not separate TF-IDF from embeddings (difference $\approx$ one standard deviation),
and it ignores cluster sizes: TF-IDF puts 18.6\% of messages in one generic cluster (``pass ride
guys thanks please''), whereas embedding clusters are balanced and name concrete problems
(accidents, waiting and cancellations, cold food orders, iOS app crashes, surge pricing). LDA fails
on short texts. \todo{human word-intrusion test to settle the comparison}.

\subsection{Probable causes (controlled simulation)}
For an incident, each metadata value (OS, app version, zone, vehicle, payment, time of day) gets a
one-sided Fisher exact test of over-representation against baseline traffic, Benjamini--Hochberg
correction, and a lift ranking. Real data has no metadata with known causes, so 1--2 causes are
planted in simulated tickets (200 trials per condition; \res{rootcause\_simulation.json}).
With 50\% purity, precision@$k$ is 0.66 / 0.86 / 0.98 for incidents of 20 / 50 / 100 tickets, against
0.35 / 0.42 / 0.33 for ranking the most frequent values. With 80\% purity both rankings exceed 0.92 (the
frequency ranking is slightly better on small incidents, where the significance requirement costs
power). On incidents without any cause, only 1.5\% report one, whereas the frequency ranking always
reports something.

\subsection{Guaranteed automation (RQ2)}
\label{sec:gate}
Learn-then-Test \citep{angelopoulos2021ltt}: on a calibration set, each of 50 pre-registered
confidence thresholds is tested with an exact binomial test of ``error among automated $> \alpha$''
at level $\delta/50$; the loosest passing threshold is kept, so that the error among automated
messages is $\leq \alpha$ with probability $\geq 1-\delta$. Fixed-sequence testing was rejected: at
the strictest thresholds too few messages are accepted to reject anything.

\begin{table}[h]
\centering\small
\begin{tabular}{lccc}
\toprule
Method & Automated & Error among automated (mean / max) & Violations \\
\midrule
Gate, $\alpha=2\%$ & 0\% (refused 10/10) & -- & 0 \\
Gate, $\alpha=5\%$ & 0\% (refused 10/10) & -- & 0 \\
\textbf{Gate, $\alpha=10\%$} & \textbf{36.8\%} & \textbf{3.9\% / 6.3\%} & \textbf{0} \\
Fixed confidence $\geq 0.8$ & 31.0\% & 3.2\% / 4.5\% & no guarantee \\
Fixed confidence $\geq 0.9$ & 8.7\% & 2.0\% / 5.3\% & no guarantee \\
\bottomrule
\end{tabular}
\caption{Complaint vs.\ not, 10 seeds, $\approx$690 calibration messages, $\delta=10\%$. Source:
\res{gate\_complaints2019.json}.}
\end{table}
A guarantee costs evidence: with this calibration size and model, $\leq 5\%$ cannot be proven, and
the gate refuses rather than pretends. \todo{automation rate vs.\ calibration size; with the
fine-tuned encoder and the gold sets}.

\subsection{Grounded replies (RQ3, generation part)}
\todo{faithfulness, citation precision/recall, policy violations, blind human comparison with real
agent replies ($\kappa$), LLM-judge validation, variance over 3 runs, cost and latency}.

\section{Error analysis and limitations}
Observed failure modes so far, each with a plausible cause:
(1) \emph{Cross-lingual gap}: without the LLM, the English-trained classical model labels the French
complaint ``J'ai été facturé deux fois\dots'' as a request (complaint probability 0.28): its training
data is English only (motivates RQ1).
(2) \emph{Verifier blind spot}: number matching accepts a wrong number that happens to appear in a
cited source (``charged 2 times'' makes ``2 days'' pass); claim-level entailment checking is needed.
(3) \emph{Retrieval}: BM25 misses paraphrases (42 of 70 questions not ranked first).
(4) \emph{Incidents}: spelling variants (\#krispykreme vs.\ ``krispy kreme'') form separate incidents.
\todo{error categories with examples for each classifier on the gold sets: sarcasm, multi-issue,
too short, label ambiguity, translationese}.

\paragraph{Threats to validity.} The Twitter data is English, from 2017, and collected around company
replies; the French evaluation relies on our own gold set; the Help Center and the customer records
are fictional, so reply-level results measure faithfulness to \emph{our} policies.

\section{Ethical considerations}
The demo company is fictional and labelled as such on every page and email; no real brand is
impersonated. Personal data is masked before storage and before any external call; demo data is
deleted after the defence. Money is never granted automatically; safety, injury, harassment and
legal cases always reach a human. Licences are respected (non-commercial, share-alike, no
redistribution of the raw data). A risk of misuse (using triage to bury complaints) is mitigated by
ranking severe and risky cases first.

\section{Conclusion and future work}
\todo{after RQ1 and RQ3 generation results}. Next steps: fine-tuned multilingual encoder, held-out
retrieval set and reranker, claim-level verifier, calibration-size study for the gate, deployment.

\section*{Contribution statement}
\todo{per member, consistent with the commit history}.

\section*{AI-use disclosure}
\todo{table: tool, purpose, extent}. An AI coding assistant was used to help write code, tests and
documentation; all numbers were produced by the scripts in the repository and checked by the team.

\bibliographystyle{plainnat}
\bibliography{references}
\end{document}
